% А1.3 — Многочлены: теорема Безу, подбор корня, деление уголком, схема Горнера
% Старая методичка: стр. 150–151 (старое задание 15)

\begin{tema}{А1.3}
\opr{Теорема Безу (следствие).} Если $P(a)=0$, то многочлен делится на $x-a$:
$P(x)=(x-a)\,Q(x)$, где степень $Q$ на единицу меньше. Так понижают степень
уравнения или неравенства.

\opr{Подбор корня.} Если коэффициенты целые, то целый корень ищут
\emph{среди делителей свободного члена}.

\begin{primer}
$x^3-2x^2+5x-24=0$. Делители 24: $\pm1,\pm2,\pm3,\dots$ Подходит $x=3$:
$27-18+15-24=0$.
\end{primer}

\begin{minipage}[t]{0.5\linewidth}
\opr{Деление уголком}\par\smallskip
\small
$\begin{array}{@{}r@{}r@{}r@{}r|l@{}}
x^3 & {}-2x^2 & {}+5x & {}-24 & \;x-3\\ \cline{5-5}
-(x^3 & {}-3x^2) & & & \;x^2+x+8\\ \cline{1-2}
 & x^2 & {}+5x & & \\
 & -(x^2 & {}-3x) & & \\ \cline{2-3}
 & & 8x & {}-24 & \\
 & & -(8x & {}-24) & \\ \cline{3-4}
 & & & 0 &
\end{array}$
\end{minipage}\hfill
\begin{minipage}[t]{0.46\linewidth}
\opr{Схема Горнера}\par\smallskip
\small
\pfGorner{3}{1 & -2 & 5 & -24}{1 & 1 & 8}{0}\par\smallskip
\footnotesize\raggedright Каждое число: предыдущее $\times\,3$ + коэффициент сверху.
Последнее — остаток.
\end{minipage}

\medskip
Итог: $x^3-2x^2+5x-24=(x-3)(x^2+x+8)$. У $x^2+x+8$ дискриминант
$D=1-32<0$, поэтому $x=3$ — единственный корень.
\end{tema}
